\subsection{Inverses with Logarithms} 
We can use the logarithm to take the inverse of any log function.
\begin{Example}
If $f(x)=2^{x-3}+5$, write an expression for $f^{-1}(x)$.
\begin{lpic}{}{5}
\regtextleft{1}{-2}{Isolate the exponential};
\regtextleft{1}{-3}{Rewrite in log form};
\regtextleft{1}{-4}{Solve for $y$};
\regtextleft{1}{-5}{Rewrite in terms of $f\inv$};
\end{lpic}
\end{Example}
The same method, in reverse, can be used to find the inverse of any logarithm function.
\begin{Example}
If $f(x)=\log_7(3x-5)+2$, write an expression for $f^{-1}(x)$.
\begin{lpic}{}{6}
\regtextleft{1}{-2}{Isolate the log};
\regtextleft{1}{-3}{Rewrite in exponential form};
\regtextleft{1}{-4}{Solve for $y$};
\regtextleft{1}{-6}{Rewrite in terms of $f\inv$};
\end{lpic}
\end{Example}